\documentclass[a4paper,11pt]{article}
\usepackage[utf8]{inputenc}
\usepackage[T1]{fontenc}
%\usepackage{lmodern}  % Clean sans-serif font
\usepackage[svgnames]{xcolor}  
\usepackage{hyperref}
\hypersetup{
    pdfborder={0 0 0}, 
    colorlinks=true,
    %linkcolor=RosyBrown,
    %urlcolor=RosyBrown
    linkcolor=gray,
    urlcolor=black
}
\usepackage{enumitem}
\usepackage[left=0.7in,right=0.7in,top=0.7in,bottom=0.7in]{geometry}
\usepackage{titlesec}
\titleformat{\section}{\large\bfseries\color{black}}{}{0em}{}[\titlerule]

\setlist[itemize,1]{left=0em, labelwidth=3em, labelsep=1em, align=left, itemsep=4pt}
\setlist[itemize,2]{left=2em, labelwidth=0em, labelsep=0.5em, align=left, itemsep=2pt}

\begin{document}

% Header
\begin{center}
    {\Large\textbf{Sofiya Zbaranska}} \\ \vspace{1.5pt}
    \textit{Physiology Ph.D. Student} \\ 
    University of Toronto \& The Hospital for Sick Children \\ 
    686 Bay St., Toronto, ON M5G 0A4 \\ \vspace{5pt}
    \href{mailto:sofiya.zbaranska@mail.utoronto.ca}{Email} $\cdot$ 
    \href{https://neur1s.github.io}{Website} $\cdot$ 
    \href{https://scholar.google.com/citations?user=7ZBuO4wAAAAJ&hl=en}{Google Scholar} $\cdot$  
    \href{https://github.com/neur1s}{Github} \\ \vspace{5pt}
    %\rule{\textwidth}{0.5pt}  % Thin dividing line
\end{center}

\section*{Summary}
I am a third-year Ph.D. student co-supervised by Drs. Josselyn and Laschowski at the University of Toronto. I am strongly committed to advancing our understanding of memory function in natural and artificial intelligence systems and display active leadership in science outreach and community engagement.  I am passionate about bridging the gap between cutting-edge research and the public, while also advocating for Ukraine and supporting those affected by the war.

% Education
\section*{Education}
\begin{description}[labelwidth=5em, labelsep=1em, leftmargin=1em, itemsep=8pt]
    \item[\textbf{2023/09 -- ongoing}] \textbf{Ph.D. in Physiology} \\ 
    \textit{Physiology Department, Temerty Faculty of Medicine, University of Toronto, Toronto, Canada}  
    \begin{itemize}
        \item Ph.D. thesis project on engram ensembles supporting social memory, supervised by Dr. Sheena Josselyn at SickKids Research Institute
        \item Research project on brain-inspired memory mechanism for neural networks based on emotional valence, supervised by Dr. Brokoslaw Laschowski at University Health Network
        \item Project on engram neurons involved in social empathy, a collaboration with Dr. Armen Saghatelyan at University of Ottawa
        \item \textbf{Average Grade:} A+
    \end{itemize}

    \item[\textbf{2021/10 -- 2023/09}] \textbf{M.Sc. in Molecular Biosciences (Neuroscience major)} \\ 
    \textit{Faculty of Biosciences, Heidelberg University, Heidelberg, Germany}  
    \begin{itemize}
        \item Master's thesis on newly identified neuronal projections involved in attention modulation, under the supervision of Dr. Hannah Monyer
        \item \textbf{Average Grade:} 1.0 (98-100\%)
        \item Student representative of the Neuroscience major
    \end{itemize}

    \item[\textbf{2018/10 -- 2021/09}] \textbf{B.Sc. in Biology} \\
    \textit{Faculty of Natural Sciences I, Martin Luther University Halle-Wittenberg, Germany}  
    \begin{itemize}
        \item Thesis on differential odorant receptor expression in male and female desert locusts, supervised by Dr. Krieger
        \item \textbf{Average Grade:} 1.5 ($\geq$ 90\%)
    \end{itemize}
\end{description}

% Honours and Awards
\section*{Honours and Awards}
\begin{itemize}
    \item[\textbf{2025/09 -- 2027/09}] Restracomp PhD Fellowship at the Hospital for Sick Children (CAD 35,000/year, competitive process)
    \item[\textbf{2025/05}] Frontiers in Physiology Best Poster Presentation Award (certificate, competitive process)
    \item[\textbf{2024/07}] BrainAid/IZN Best Thesis Award (certificate and book prize, competitive process)
    \item[\textbf{2022/10 -- 2023/10}] Graduate Scholarship from the Heinrich Böll Foundation for studies at Heidelberg University (EUR 10,892/CAD 16,000, competitive process)
    \item[\textbf{2022/8 -- 2022/11}] Master's Fellowship at Karolinska Institutet (SEK 48,000/CAD 6,000, \\ non-competitive, through Karolinska Institute research grants)
    \item[\textbf{2019/10 -- 2020/10}] German Scholarship for Talented and Oustanding Students at Martin Luther University Halle-Wittenberg (EUR 3,600/CAD 5,400, competitive process)
    \item[\textbf{2015/05}] Honours List with Distinction, High School No. 329 `Logos', Kyiv, Ukraine 
    \item[\textbf{2014/06}] Second Place in the District Level English Olympiad (High School No. 329 `Logos', Kyiv, Ukraine, competitive process)
\end{itemize}

% Publications
\section*{Publications}
\subsection*{Peer-reviewed}
\begin{enumerate}
    \item Park, S., Wu, B., \textbf{Zbaranska, S.}, Lee J., Jacob A.D., Hoorn A., Mocle A., Luchetti A., Jung J.H., Mocle A.J.,  Frankland P.W., Josselyn S.A. (2026) Temporally-gated offline engram ensemble reverberation in the lateral amygdala is required for fear memory consolidation. Accepted at \textit{Proceedings of the National Academy of Sciences}.\\
    \textbf{Role:} Analyzed behavioural data and applied machine learning techniques to calcium imaging datasets; contributed to manuscript writing and revision.
    \textbf{Contribution:} 11--20\%

    \item \textbf{Zbaranska, S.}, Josselyn, S. A. (2025). Astrocytes: Emerging Stars of Engrams. \textit{Cell Research} (IF: 28.2) 35(4), 241-242. \href{https://doi.org/10.1038/s41422-024-01066-4}{https://doi.org/10.1038/s41422-024-01066-4} (Refereed). \\
    \textbf{Role:} First author. Performed literature search, wrote the manuscript, prepared the figure, and participated in editing and revisions.  
    \textbf{Contribution:} 90\%

    \item \textbf{Zbaranska, S.}, Lee, J.J., Zaheer, M., Josselyn S.A. (2024). Unraveling Engrams: Tracing Memories through Time. \textit{Reference Module in Neuroscience and Biobehavioral Psychology}. \\
    \href{https://doi.org/10.1016/b978-0-443-15754-7.00040-7}{https://doi.org/10.1016/b978-0-443-15754-7.00040-7} (Refereed). \\
    \textbf{Role:} Co-first author. Performed literature search, wrote 1/3 of the manuscript, prepared figures, and participated in editing and revisions.  
    \textbf{Contribution:} 31--40\%
\end{enumerate}

\subsection*{In Preparation}
\begin{enumerate}
    \item \textbf{Zbaranska, S.}, Rajeev A., Josselyn S., Laschowski B. (2026) A Brain-inspired Memory Mechanism for Neural Networks Based on Emotional Valence (under preparation for \textit{IEEE Transactions on Neural Networks and Learning Systems}). \\
    \textbf{Role:} Project lead, designed the model architecture, specified the tasks and training procedure, wrote the code, performed data analysis, and wrote the manuscript.
    \textbf{Contribution:} 70\%

    \item Zanjanpour D., \textbf{Zbaranska, S.}, Liu H.T., Laschowski B. A Neuroscientific Framework for Trust in Human-Robot Interaction. \\
    \textbf{Role:} Literature review, conceptual framework development, and manuscript writing.
    \textbf{Contribution:} 20\%

    \item Luchetti, A., Wu, B., Hoorn A.A., \textbf{Zbaranska, S.}, Jacob A.D., Mocle A.J,  Chen M., DeCristofaro A., Frankland P.W., Josselyn S.A. Decoding of Neuronal Activity Patterns in the Lateral Amygdala during Fear Memory Formation in Mice. \\
    \textbf{Role:} Analyzed behavioural data and analyzed the role of different neuronal populations using machine learning classifiers. 
    \textbf{Contribution:} 10\%

\end{enumerate}

% Journal Review Activities
\section*{Peer Review Activities}
\begin{itemize}
    \item[\textbf{2026/02 -- 2026/03}] \textit{International Conference on Learning Representations}: reviewed 2 manuscripts.

    \item[\textbf{2023/12}] \textit{Philosophical Transactions of the Royal Society B}: reviewed 1 manuscript.
\end{itemize}

% Journal Review Activities
\section*{Public Talks}
\begin{itemize}
    \item[\textbf{2026/02/05}] \textit{Modeling Memory Processes in the Brain Using Transformers} (The Hospital for Sick Children, Toronto, ON Canada). Invited by AI4SickKids Reading Group orginizing committee.
    \item[\textbf{2025/11/27}] \textit{Telling Friends from Foes: Contributions of the Medial Amygdala to Social Memory} (The Hospital for Sick Children, Toronto, ON Canada). Invited by the Neuroscience and Mental Health Program Council at the Hospital for Sick Children.
    \item[\textbf{2025/10/13}] \textit{The Factory of Memories: Taming the Engram} (watch on \href{https://youtu.be/ETVgqyQwZVM?si=o-6VG4Vc8KOKGTGn}{\textbf{Youtube}}). Invited by the Faculty of Applied Sciences at Ukrainian Catholic University.
\end{itemize}

\section*{Conference Presentations}
\begin{enumerate}
    \item \textbf{Zbaranska S.}, Sarikahya M.H., Luchetti A., Frankland P.W., Josselyn S.A. How the Medial Amygdala Shapes Neutral and Emotionally Charged Social Memories. 2026 Annual Meeting of the Society for Social Neuroscience (upcoming). Montreal, QC, Canada, August 2026. Poster presentation.
    \item \textbf{Zbaranska S.}, Rajeev A., Josselyn S., and Laschowski B. Neuroscience-inspired AI for memory. CPIN Research Day 2026 (upcoming). Toronto, ON, Canada, June 2026. Poster presentation.
    \item \textbf{Zbaranska S.}, Sarikahya M.H., Luchetti A., Frankland P.W., Josselyn S.A. Telling Friends from Foes: Contributions of the Medial Amygdala to Social Memory. Canadian Association for Neuroscience Meeting 2026 (upcoming). Montreal, QC, Canada, May 2026. Poster presentation.
    \item \textbf{Zbaranska S.}, Frankland P.W., Josselyn S.A. Telling Friend from Foe: The Role of the Medial Amygdala in Social Memory. CPIN Research Day 2025. Toronto, ON, Canada, June 2025. Poster presentation.
    \item \textbf{Zbaranska S.}, Frankland P.W., Josselyn S.A. Contributions of the Medial Amygdala to Social Memory. Annual Frontiers in Physiology 2025. Toronto, ON, Canada, May 2025. Poster presentation. \textit{Recipient of the Best Poster Presentation Award}.
    \item \textbf{Zbaranska S.}, Frankland P.W., Josselyn S.A. Identifying Engrams Supporting Social Memory in the Medial Amygdala. Canadian Association for Neuroscience Meeting 2025. Toronto, ON, Canada, May 2025. Poster presentation.
    \item \textbf{Zbaranska S.}, Held K., Chowdhury D., MacLaren D.A.A., Monyer H. Characterization of Medial Septal SOM-Positive Neurons and Their Projections to the Hippocampus. IZN Retreat 2023. Closter Schoental, BW, Germany, July 2023. Poster presentation.
\end{enumerate}

% Research Experience
\section*{Research Experience}
\begin{itemize}
    \item[\textbf{08/2025 -- ongoing}] \textbf{Collaborative Computational Research Project with Dr. Laschowski} \\
    \textit{University Health Network, Toronto, Canada} \\
    Designed and implemented a transformer model trained on a synthetic episodic memory task. Incorporated emotional valence levels to incentivize memory prioritization. Showed that this approach improves long-term memory performance, spatial information representation, and robustness to noise. Manuscript is in preparation for submission to \textit{IEEE Transactions on Neural Networks and Learning Systems}.

    \item[\textbf{09/2023 -- ongoing}] \textbf{Ph.D. Research Project in Josselyn lab} \\
    \textit{The Hospital for Sick Children, Toronto, Canada} \\
    Designed and am leading a project focusing on the role of the medial amygdala in social memory. An additional branch of my project addresses how engrams in the prefrontal cortex support empathy in mice (collaboration with Dr. Armen Saghatelyan at University of Ottawa). My work includes mouse neurosurgeries, genetic tools, behavioural assays, in vivo calcium imaging, and advanced data analysis.

    \item[\textbf{03/2023 -- 09/2023}] \textbf{Master's Thesis in Monyer lab} \\
    \textit{German Cancer Research Center, Heidelberg, Germany} \\
    Performed electrophysiological recordings in freely behaving mice to study the role of a newly dicovered class of projections from the medial septum to the hippocampus (by testing their involvement in attention modulation).

    \item[\textbf{12/2022 -- 02/2023}] \textbf{Lab Rotation in Liu lab} \\
    \textit{German Cancer Research Center, Heidelberg, Germany} \\
    Studied the role of the chromatin remodeller CHD7 in the hippocampal neurogenesis
    using immunohistochemistry and single-cell RNA sequencing.

    \item[\textbf{08/2022 -- 11/2022}] \textbf{Research Internship in Ehrsson lab} \\
    \textit{Karolinska Institute, Stockholm, Sweden} \\
    Conducted a research project about the link between body ownership and formation of
    false memories using virtual reality in humans. This encompassed behavioural experiments with human volunteers, data collection, and analysis.

    \item[\textbf{05/2022 -- 07/2022}] \textbf{Lab Rotation in Kuner lab} \\
    \textit{Interdisciplinary Center for Neuroscience, Heidelberg University, Germany} \\
    Investigated the influence of isoflurane anesthesia on the cerebral cortex using in vivo two-photon imaging. Gained experience in mouse neurosurgeries (viral injections and cranial window implantations), data collection and analysis of longitudinal 2-photon fluorescence data.

    \item[\textbf{02/2022 -- 04/2022}] \textbf{Lab Rotation in Monyer lab} \\
    \textit{German Cancer Research Center, Heidelberg, Germany} \\
    Examined the mechanisms of environment representation and association formation using in vivo electrophysiological recordings in the lateral entorhinal cortex in freely moving mice. Contributed to the publication by \href{https://doi.org/10.1016/j.neuron.2023.06.020}{Huang et al. (2023) \textit{Neuron}}.

    \item[\textbf{10/2021 -- 01/2022}] \textbf{Lab Rotation in Bading lab} \\
    \textit{Interdisciplinary Center for Neuroscience, Heidelberg University, Germany} \\
    Performed whole-cell patch clamp recordings of humanized mouse hippocampal neurons, analyzed the acquired data and contributed to the publication by \href{https://doi.org/10.1016/j.jbc.2023.104671}{Pruunsild et al. (2023) \textit{J Biol Chem.}} 

    \item[\textbf{03/2021 -- 09/2021}] \textbf{Bachelor's Thesis project in Krieger lab} \\
    \textit{Department of Zoology, Martin Luther University Halle-Wittenberg, Germany} \\
    Extracted mRNA from desert locust antennae, performed reverse transcription, and established a quantitative real-time PCR protocol to assess sex- and stage-specific expression of odorant receptors.

    \item[\textbf{06/2020 -- 08/2021}] \textbf{Undergraduate Research Assistant in Bonas lab} \\
    \textit{Department of Plant Genetics, Martin Luther University Halle-Wittenberg, Germany} \\
    Research focus on plant pathogenic bacteria.
    Tasks: cloning (Golden Gate method), transformation, plasmid and RNA purification, restriction analysis, PCR, gel electrophoresis, SDS-PAGE and Western Blot.
\end{itemize}

% Leadership and Engagement
\section*{Leadership and Extracurricular Activities}
\subsection*{Within the University}
\begin{itemize}
    \item[\textbf{2026/02-2026/09}] \textbf{AI4SickKids Reading Group Organizer}, The Hospital for Sick Children \\ 
    Contributing to the organization of the AI4SickKids Reading Group aimed at disseminating AI- and ML-related research to SickKids community. Involved in speaker recruitment, scheduling, and event promotion. Delivered a talk on modelling memory in the brain using transformers.
    \item[\textbf{2026/01-ongoing}] \textbf{Neurosciences \& Mental Health Trainee Council Member}, The Hospital for Sick Children \\ 
    Participating in the organization of the Neurosciences \& Mental Health Research Day 2026 by contrubuting to speaker recruitment, planning of trainee-focused events such as workshops and social gatherings.
    \item[\textbf{2025/02-2025/05}] \textbf{Search Committee Member}, The Hospital for Sick Children \\ 
    Represented graduate student perspective on the search committee for recruiting a Scientist/Senior Scientist in Molecular and Cellular Basis of Cognition. Contributed to application review, candidate interviews, and scoring.
    \item[\textbf{2025/05 and 2024/05}] \textbf{Volunteer}, Science Rendezvous, St. George Campus, University of Toronto \\  
    Coordinated poster presentations, engaged with the public and facilitated event logistics.
    \item[\textbf{2024/07-2025/03}] \textbf{Mentor for Undergraduate Students}, The Hospital for Sick Children \\ 
    Mentored two undergraduate students in Dr. Josselyn's lab: one received hands-on training in wet lab techniques, and the other was guided through machine learning-based classification of neuronal types in calcium imaging data.
    \item[\textbf{2025/07 and 2024/05}] \textbf{Kids Science Lab Tours}, The Hospital for Sick Children \\ 
    Held lab tours for high school students and International Brain Bee competition awardees, demonstrated microscopy, promoted our research, and explained the role of machine learning in neuroscience.
    \item[\textbf{2021/10-2023/08}] \textbf{Neuroscience Major Representative}, Heidelberg University \\ 
    Facilitated communication between students and faculty in the Molecular Biosciences Master's program. Addressed student inquiries on program planning, lab rotations, and international research internships. Organized and led information sessions for prospective students. Coordinated social events (e.g., pub quizzes, hikes) to strengthen student community.
    \item[\textbf{2021/12-2022/10}] \textbf{Symposium Co-organizer}, Heidelberg University \\ 
    Recruited speakers and helped run the “Sustainability in Science” symposium. Moderated sessions and introduced speakers. Organized social events such as a pub quiz and group hike.
\end{itemize}

\subsection*{Outside the University}
\begin{itemize}
    \item[\textbf{2023/09 -- ongoing}] \textbf{Active Member}, Ukrainian Students' Club in Toronto \\ 
    Organizing cultural and fundraising events to support the Ukrainian community in Toronto, such as podium discussions with the Consul General of Ukraine.
    \item[\textbf{2022/03 -- 2022/08}] \textbf{Volunteer}, Heidelberg \\ 
    Served as a translator for Ukrainian refugees for doctor visits, temporary permit applications, and bank appointments.
    \item[\textbf{2019/06 -- 2021/08}] \textbf{High School Student Mentor}, online \\ 
    Advised Ukrainian high school graduates on applying to bachelor's and master's programs in Germany. Supported university applications, housing searches, bank account setup, and understanding of German health insurance.

\end{itemize}

% Conference Participation
\section*{Conference Participation}
\begin{itemize}
    \item[\textbf{2025/09}] Neuronto 2025, Toronto ON, Canada 
    \item[\textbf{2025/06}] Collaborative Program in Neuroscience (CPIN) Research Day 2025, Toronto ON, Canada 
    \item[\textbf{2025/05}] Annual Frontiers in Physiology Research Symposium 2025, Toronto ON, Canada 
    \item[\textbf{2025/05}] Canadian Association for Neuroscience Meeting 2025, Toronto ON, Canada 
    \item[\textbf{2024/11}] Neuroscience and Mental Health Research Day 2024, SickKids, Toronto ON, Canada 
    \item[\textbf{2024/06}] Neuronto Research Symposium, MaRS Centre, Toronto ON, Canada 
    \item[\textbf{2023/11}] Neuroscience and Mental Health Trainee Conference, SickKids, Toronto ON, Canada 
    \item[\textbf{2023/07}] Interdisciplinary Center for Neurosciences Retreat 2023, Schoental BW, Germany 
    \item[\textbf{2022/09}] Symposium "Sustainability in Science", Heidelberg BW, Germany
    \item[\textbf{2021/09}] Symposium "Science Outside the Box", Heidelberg BW, Germany
\end{itemize}

% Skills and Technical Expertise
\section*{Certifications}
\begin{itemize}
    \item[\textbf{2025/07}] NeuroAI (theory and project components), \textit{Neuromatch}
    \item[\textbf{2025/07}] Computational Neuroscience (self-paced), \textit{Neuromatch}
    \item[\textbf{2023/01}] Course on Laboratory Animals - Function A, D and C \\
    \textit{Federation of European Laboratory Animal Science Associations, Germany}
    \item[\textbf{2022/09}] MATLAB Fundamentals, \textit{MathWorks}
    \item[\textbf{2022/03}] Intermediate R, \textit{DataCamp}
    \item[\textbf{2022/01}] Research Integrity Course, \textit{Epigeum}
    \item[\textbf{2022/01}] Biology Meets Programming (Bioinformatics), \textit{UC San Diego}
\end{itemize}


\end{document}
